\documentclass[12pt]{article}

\usepackage[utf8]{inputenc}
\usepackage[T1]{fontenc}
\usepackage[brazil]{babel}

\usepackage{geometry}
\geometry{margin=2.5cm}
\usepackage{setspace}
\onehalfspacing

\usepackage{amsmath, amssymb, amsthm}
\usepackage{mathtools}

\numberwithin{equation}{section}

\theoremstyle{plain}
\newtheorem{theorem}{Teorema}[section]
\newtheorem{proposition}[theorem]{Proposição}
\newtheorem{lemma}[theorem]{Lema}
\newtheorem{corollary}[theorem]{Corolário}

\theoremstyle{definition}
\newtheorem{definition}[theorem]{Definição}
\newtheorem{example}[theorem]{Exemplo}
\newtheorem{exercise}{Exercício}[section]

\theoremstyle{remark}
\newtheorem*{remark}{Observação}

\usepackage{hyperref}
\usepackage{csquotes}

\usepackage[
    backend=biber,
    style=authoryear,
    maxcitenames=2
]{biblatex}

\addbibresource{/mnt/c/Users/nicho/Documents/nicholasvoltani.github.io/bibliography.bib}

\newcommand{\R}{\mathbb{R}}
\newcommand{\pref}{\succeq}
\newcommand{\strictpref}{\succ}
\newcommand{\indiff}{\sim}

% ------------------------------------------------
% Title (fill manually)
% ------------------------------------------------
\title{}
\author{}
\date{}

\begin{document}

\maketitle

\begin{longtable}[]{@{}l@{}}
\toprule\noalign{}
\endhead
\bottomrule\noalign{}
\endlastfoot
date: ``2026-03-30'' \\
tags: \\
- economics \\
aliases: \\
\end{longtable}

\hypertarget{conceitos-introdutuxf3rios}{%
\section{Conceitos introdutórios}\label{conceitos-introdutuxf3rios}}

\hypertarget{aversuxe3o-ao-risco}{%
\subsection{Aversão ao risco}\label{aversuxe3o-ao-risco}}

\begin{definition}
O equivalente certo de um tomador de risco é a quantidade monetária
não-aleatória com a qual ele obtém utilidade igual à sua
\textit{utilidade média}.

Dada uma loteria \(L\) e uma utilidade de Bernoulli \(u(\cdot)\), o
equivalente certo (ou equivalente de certeza) \(C(L, u)\) é o valor
(monetário) \textit{com probabilidade \(100\%\) de ser recebido} que faz
com que a utilidade de Bernoulli se iguale à Utilidade de von
Neumann-Morgenstern \(U(\cdot)\) da loteria \(L\),
i.e.~\(U(L)\).\footnote{\textcite{Cowell2004}, p.~188 comenta sobre a
  confusão que ambas estas utilidades --- utilidade ``de Bernoulli''
  (cf. \textcite{Mas-Colell1995}, p.~184) e utilidade de von
  Neumann-Morgenstern --- podem suscitar: ``Here we encounter a
  terminologically awkward corner. We should not really call \(u\)
  {[}''utilidade de Bernoulli''{]} `the utility function' because the
  whole expression {[}utilidade de von Neumann-Morgenstern,
  \(U(L) = \sum \limits_{k} p_{k} u_{k}\){]} is the person's utility; so
  \(u\) is sometimes known as the individual's \textit{cardinal utility
  function} or \textit{felicity function}; arguably neither term is a
  particularly happy choice of words.''.}

Formalmente, temos que, para uma loteria \((p_{1}, \dots, p_{n})\)
relacionada às utilidades (de Bernoulli)
\((u_{1} \coloneqq u(x_{1}), \dots, u_{n}:= u(x_{n}))\), a definição do
equivalente certo segue a forma
\[
u(C(L, u)) = \sum \limits_{k} p_{k} u_{k} = U(L)
\]

Pensando em termos de gráficos: \(c\) é o valor (não-aleatório,
``\textit{certo}'') com o qual o indivíduo alcança a utilidade média
(dadas as suas possibilidades \(u_{1},\dots, u_{n}\)). Intuitivamente,
pensaríamos que este ponto seria
\[
c = \sum \limits_{i=1}^{n} p_{i} x_{i}
\]
i.e.~o valor esperado dentre as possibilidades \(x_{i}\). Este caso,
porém, é o que se chama de ``indiferente ao risco''. Indivíduos que têm
\textit{aversão} ao risco, por outro lado, preferem um valor
\textit{menor, porém certo}, para evitar a aleatoriedade; eles se
dispõem a receber um valor menor, contanto que não se indisponham a
``recebíveis'' incertos. \textit{Mutatis mutandis} para indivíduos
\textit{propensos} ao risco: eles precisam ser pagos valores certos
\textit{maiores} para que se abstenham do risco.

!445 Fonte: \textcite{Cowell2004}, p.~191. Na figura, os possíveis
acontecimentos são \(x_{\color{red} Red}, x_{\color{blue} Blue}\)
(implicitamente com probabilidades respectivas
\(p_{\color{red} Red}, p_{\color{blue} Blue}\)), \(\mathcal{E}x\) é o
valor esperado dos acontecimentos \(x\) (i.e.~\(\sum p_{i}x_{i}\)), e
\(\xi\) é o equivalente certo para este indivíduo. Note-se que, como
\(\xi < \mathcal{E}x\), este indivíduo é \textit{avesso} ao risco; isso
também é notável por sua utilidade \textit{côncava}
(\(u^{\prime\prime}<0\)).
\end{definition}

\begin{definition}
Dado um indivíduo com um equivalente certo \(C(L, u)\) e cujos possíveis
acontecimentos \(x_{1},\dots,x_{n}\) (com resp. probabillidades
\(p_{1},\dots,p_{n}\)) tenham valor esperado
\[
\left<x\right> \coloneqq \sum \limits_{i=1}^{n} p_{i} x_{i}
\]
define-se seu \textit{prêmio de risco} como
\(\left< x \right> - C(L,u)\)
\parencite[, p. 191]{Cowell2004}.\footnote{Note-se que esta definição
  \textbf{não} é a mesma que \textcite{Mas-Colell1995} faz: MWG definem
  o que chamaram de ``\textit{probability premium}'' (definição
  \textbf{6.C.2 (ii)}, p.~186), mas não o \textit{risk premium}
  \textit{per se}.}

Para indivíduos indiferentes ao risco, ele é \(= 0\); para os avessos ao
risco, é \(>0\); para os propensos ao risco, é \(<0\).
\end{definition}

\hypertarget{exercuxedcios}{%
\section{Exercícios}\label{exercuxedcios}}

\begin{exercise}
\parencite[, pp. 233-5]{Varian2014} Suponhamos que um indivíduo possui
renda \(w\) e queira investir uma parcela \(x\) dela em algum ativo
arriscado, tal que divida sua renda de forma
\[
\begin{cases}
x&: \text{Ativo arriscado}\\
w-x&: \text{Riqueza não-investida} 
\end{cases}
\]

Este ativo pode ter um retorno \(r_{b}\) em um cenário ``ruim'' com
probabilidade \(p\), e um retorno \(r_{g}\) em um cenário ``bom'' com
probabilidade \(1-p\).\footnote{Caso vá checar no Varian, ele utiliza
  \(p\) para o resultado \textit{bom} e \(1-p\) para o resultado
  \textit{ruim}. Caso haja algum erro aqui, com certeza se deu devido a
  esta alteração aqui.} Consideremos, em geral, que \(r_{g}> 0\), e que,
ao menos, \(r_{b}<r_{g}\), em geral tendo \(r_{b}<0 < r_{g}\).

A riqueza deste indivíduo, em ambos os cenário, será

\begin{align*}
W_{g} = (w-x) + x (1+r_{g}) = w + x r_{g}\\
W_{b} = (w-x) + x (1+r_{b}) = w + x r_{b}
\end{align*}

A utilidade esperada deste indivíduo será
\[
\left<u\right> = p \,u(W_{b}) + (1-p) \,u(W_{g})
\]
e este indivíduo deseja encontrar \(x\) ótimo que maximize sua utilidade
esperada. Supomos também que ele é avesso ao risco.

No que tange à derivada da utilidade esperada com relação a \(x\),
investigue o caso em que \(x=0\). Qual é a condição com que se pode
garantir que haja \(x^*>0\)? Como ela relaciona-se com
\(\langle r \rangle = p r_{b}+(1-p)r_{g}\)?

Satisfeitas as condições necessárias para que \(x^*>0\), investigue o
que ocorre quando temos uma taxação sobre os retornos:
\[
r_{i} \mapsto \tilde{r}_{i} = (1-t) r_{i}
\]

Mostre que a solução ótima de riqueza pós-taxação \(\hat{x}\) de se
investir satisfaz
\[
\hat{x} = \frac{x^*}{1-t}
\]
quando \(x^*\) é solução ótima pré-taxação. Isso condiz com sua intuição
de \(\hat{x} \gtrless x^*\)?

Sinta-se livre para explorar este exercício nesta simulação:
\href{https://www.desmos.com/calculator/ppkg28wznf}{Investimentos e
Taxação (Simulação no Desmos)}!
\end{exercise}

\begin{exercise}
\parencite[, pp. 223-4]{Cowell2004} Um indivíduo precisa declarar sua
renda \(y\) às autoridades. A renda que ele pode declarar,
\(0 \leq x \leq y\), será taxada a uma taxa \(t \in (0,1)\).

Suponha que a probabilidade de que sua renda declarada seja
\textit{auditada} é \(p\). Caso o indivíduo declare \(x < y\), então ele
deverá pagar a mais como multa, de tal forma que sua renda disponível
fica na forma:
\[
\begin{cases}
\text{Sem auditoria: } y - tx \\
\text{Com auditoria: } (1-(1+s)t)y + stx 
\end{cases}
\]

Suponha que o indivíduo é avesso ao risco, e que tenha utilidade
\(u(\cdot)\). Escreva a utilidade média \(\langle u \rangle\), e obtenha
a condição de primeira ordem da maximização da utilidade média (com
relação a \(x\)).

Através da condição de primeira ordem, mostre que, quando
\(s=\frac{1-p}{p}\), vale que \(x=y\), e que quando
\(s>\frac{(1-p)}{p}\), valerá que \(x < y\).

\begin{remark}
O motivo de o segundo caso ficar dessa forma é o seguinte. Supondo que
\textit{vá haver auditoria}, temos os casos extremos 1) para \(x=y\),
naturalmente ele terá \((1-t)y\) de renda disponível em ambos os casos
(afinal, quem não deve, não teme); 2) com \(x=0\), ele fica com
\((1 - (1+s)t)y\) (o imposto pago é cobrado a mais). Por regra de
três/triângulos equivalentes, temos que

\begin{align*}
\frac{[(1-(1+s)t)y] - [(1-t)y]}{y - 0} = \frac{Y - [(1-t)y]}{(y - x) - 0} \\
\therefore Y = (1-(1+s)t)y + stx
\end{align*}
\end{remark}
\end{exercise}

\begin{exercise}
Mostre que o equivalente certo de uma função utilidade exponencial
\(u(x) = - e^{-ax}\) assume a forma
\[
C = - \frac{1}{a} \ln\left( \sum \limits_{i}p_{i}e^{-ax_{i}} \right) = \frac{1}{a} \ln\left( \frac{1}{\sum \limits_{i}p_{i}e^{-ax_{i}}} \right)
\]
Fazendo a transformação \(x_i = \ln \frac{\tilde{x_{i}}}{a}\) e
\(C = \ln \frac{\tilde{C}}{a}\), esta fórmula assume a forma de uma
\textit{média harmônica ponderada}
\[
\tilde{C} = \frac{1}{\sum \limits_{i} p_{i} \frac{1}{\tilde{x}_{i}}} 
\]

\end{exercise}

\begin{exercise}
Mostre que o equivalente certo de uma função utilidade isoelástica
\(u(x) = \frac{x^{1-\rho}-1}{1-\rho}\) assume a forma de uma
\textit{\((1-\rho)\)-norma ponderada}, i.e.~
\[
   C(L, \rho) = \left( \sum \limits_{i} p_{i} x_{i}^{1-\rho} \right)^{1/(1-\rho)}
   \]

\end{exercise}

\begin{exercise}
Mostre que o equivalente certo de uma função utilidade logarítmica
\(u(x) = \ln x\) tem a forma de uma \textit{média geométrica ponderada}
\[
C(L) = \prod \limits_{i=1}^{n} x_{i}^{p_{i}}
\]

\end{exercise}

\begin{exercise}
\parencite[, p. 230-1]{Varian2014} Suponha uma comunidade de \(N\)
indivíduos, cada um com renda \(w\) e com uma probabilidade
\(0\leq p \leq 1\) de sofrer uma perda \(L\) (ou seja, ficando com renda
\(w-L\)).

Suponha que estes indivíduos reservem uma quantidade \(q\) de suas
rendas (ficando com renda \(w-q\)) em um ``\textit{pool}'' comum, de tal
forma que indivíduos que sofram perdas possam ressarcir-se a partir
deste \textit{pool}.

Supondo que os riscos dos indivíduos sejam \textit{independentes entre
si}, calcule qual tem de ser essa quantidade \(q\) em função de \(p\) e
\(L\). Ela depende de \(N\)? Por quê/por que não?

Ademais, verifique que este valor é justamente o que um indivíduo
\textit{indiferente ao risco} pagaria por este seguro: se sua função
utilidade for \(u(x) = ax\) (\(a>0\)), então ele pagaria este mesmo
valor. (Dica: a fórmula do equivalente certo diz sobre \textit{a renda
equivalente}, enquanto aqui queremos \textit{o valor subtraído da renda}
que ele aceita para se desfazer do risco.)

Obs: Isto é o que se costuma chamar de \textit{mutualismo}, ou
\textit{risk pooling}, em ciências atuariais.

Obs²: Sinta-se à vontade para testar este exemplo com algumas funções
utilidade comuns nesta simulação:
\href{https://www.desmos.com/calculator/quce5bk5v5}{Equivalentes certos
(Simulação Desmos)}.
\end{exercise}

\hypertarget{resoluuxe7uxf5es}{%
\section{Resoluções}\label{resoluuxe7uxf5es}}

\begin{exercise}
Diferenciando \(\langle u \rangle\) com relação a \(x\), temos
\[
\frac{d\langle u \rangle}{dx} = p r_{b} \, u^\prime(W_{b}) +(1-p) r_{g} \, u^\prime(W_{g})
\]

A segunda derivada fica como
\[
\frac{d^{2}\langle u \rangle}{dx^{2}} = p r_{b}^{2} \, u^{\prime\prime}(W_{b}) +(1-p) r_{g}^{2} \, u^{\prime\prime}(W_{g}) < 0
\]
que é \(<0\) por conta da aversão ao risco pressuposta.

É pertinente que consideremos o caso da derivada de
\(\langle u \rangle\) em que \(x=0\), pois, caso sua utilidade esperada
\textit{somente caia} para \(x>0\), ele não investirá nada, pois terá
seu máximo em \(x=0\). Seu valor é

\begin{align*}
\frac{d\langle u \rangle}{dx}(x=0) &= p r_{b} u^\prime(w) + (1-p) r_{g} u^\prime \\
&= u^\prime(w) \langle r \rangle 
\end{align*} onde \(\langle r \rangle=p r_{b} + (1-p) r_{g}\) é o
retorno esperado. Naturalmente, ele estará mais instigado a investir
caso \(\langle r \rangle>0\)\footnote{Por ser avesso ao risco. Não tende
  a ser o mesmo para propenso ao risco.}. Portanto, tomando também que
\(u^\prime(w)>0\)\footnote{Pois sempre pressupomos que o indivíduo
  prefere mais dinheiro do que menos --- ou seja, derivada positiva ---,
  não obstante ele ter ``utilidade marginal decrescente''. Só porque o
  Tio Patinhas se satisfaz menos com uma moeda a mais em seu cofre cheio
  do que nele vazio, não quer dizer que ele não se satisfaça
  \textit{algum pouquinho a mais}!}, então sua utilidade marginal
tenderá a aumentar caso \(\langle r \rangle>0\), portanto instigando-o a
investir alguma parcela \(x>0\) de sua riqueza neste ativo.

Voltando ao caso geral: supomos que a função utilidade é bem-comportada
o suficiente para que encontremos \(x^*\) ótimo através de
\(\frac{d\langle u \rangle}{dx}=0\).

Buscamos agora o efeito de uma taxação sobre os rendimentos deste
retorno. Ou seja, os retornos pós-taxação serão \((1-t)r_{g}\) e
\((1-t)r_{b}\)\footnote{A rigor, não faria muito sentido taxar no caso
  \(r_{b}<0\) --- o indivíduo seria \textit{ressarcido} (parcialmente)
  pelo governo por sua perda! Desconsideremos este caso, por
  simplicidade\ldots{} Mas note que este ``deslize'' desempenha um papel
  não-desprezível no resultado obtido deste exercício.}, de forma que
sua riqueza fica da forma

\begin{align*}
W_{g} = (w-(1-t)x) + (1-t)x (1+r_{g}) = w + (1-t)x r_{g}\\
W_{b} = (w-(1-t)x) + (1-t)x (1+r_{b}) = w + (1-t)x r_{b}
\end{align*} A utilidade marginal fica da mesma forma, embora
\(W_{i}(t>0) < W_{i}(t=0)\). Sua derivada fica numa forma similar:

\begin{align*}
\frac{d\langle u \rangle}{dx} &= p (1-t)r_{b}\,u^\prime(w + (1-t)x r_{b}) + (1-p) (1-t)r_{g}\,u^\prime(w + (1-t)x r_{g})
\end{align*}

Chame-se sua solução ótima pós-taxação de \(\hat{x}\). Note-se que a
solução
\[
\hat{x} = \frac{x^*}{1-t}
\]
satisfaz que \(\frac{d\langle u \rangle}{dx}=0\):

\begin{align*}
\frac{d\langle u \rangle}{dx} &= p (1-t)r_{b}\,u^\prime\left( w + \frac{\cancel{ (1-t) }x^*}{\cancel{ 1-t }} r_{b} \right) + (1-p) (1-t)r_{g}\,u^\prime\left( w + \frac{\cancel{ (1-t) }x^*}{\cancel{ 1-t }} r_{g} \right) \\
&= 0
\end{align*} pois torna-se a condição ótima de \(x^*\) (pré-taxação).

Note-se, portanto, que aumentar \(t\) \textbf{aumenta o percentual de
riqueza investida}! A explicação de Varian é a seguinte

\begin{quote}
``Mediante imposto, o indivíduo terá um ganho menor no cenário bom,
\textit{mas também terá uma perda menor no cenário ruim}. Ao aumentar
seu investimento inicial em \(1/(1 - t)\), o consumidor pode reproduzir
os mesmos retornos \textit{pós-impostos} que obtinha antes da
implementação do imposto. O imposto reduz seu retorno esperado, mas
também reduz seu risco: ao aumentar seu investimento, o consumidor pode
obter exatamente o mesmo padrão de retornos que tinha antes e, assim,
compensar completamente o efeito do imposto. Um imposto sobre um
investimento arriscado representa um subsídio sobre a perda quando o
retorno é negativo.'' \parencite[, p. 235]{Varian2014}
\end{quote}

Relembrando: sinta-se livre para explorar este exercício nesta
simulação:
\href{https://www.desmos.com/calculator/ppkg28wznf}{Investimentos e
Taxação (Simulação no Desmos)}!
\end{exercise}

\begin{exercise}
Escrevendo \(y_{n} = y-tx\) e \(y_{a}=(1-(1+s)t)y + stx\) as rendas
disponíveis no caso sem e com auditoria, respectivamente. Então a
condição de primeira ordem fica como (já cancelando o fator \(t\) devido
ao \(=0\))

\begin{align*}
(1-p) u^\prime(y_{n}) + sp u^\prime(y_{a}) \mp sp u^\prime(y_{n}) &= 0 \\
(1-p - sp) u^\prime(y_{n}) + sp (\underbrace{ u^\prime(y_{n}) - u^\prime(y_{a}) }_{ \leq 0 }) &=0
\end{align*}

Como \(y_{a} \leq y_{n}\) --- e, evidentemente,
\(y_{a}=y_{n} \iff x = y\) ---, temos que, embora
\(u(y_{a}) \leq u(y_{n})\) (``mais é melhor''), temos que
\(u^\prime(y_{a}) \geq u^\prime(y_{n}) \iff u^\prime(y_{n})- u^\prime(y_{a}) \leq 0\),
por conta de utilidades marginais decrescentes. Porém, a condição de
primeira ordem somente é satisfeita quando, além de \(x=y\),
\[
s = \frac{1-p}{p}
\]

Caso \(1- p - sp < 0 \iff s > \frac{1-p}{p}\), a condição de primeira
ordem nunca pode ser \(=0\), pois a expressão inteira é \(<0\) --- a não
ser que o indivíduo declare, \textit{voluntariamente}, ter \textit{mais}
do que ele tem (i.e.~se dispor a pagar por renda que ele \textit{não}
tem)!

Portanto, caso \(s < \frac{1-p}{p}\), teremos que a escolha ótima do
indivíduo será reportar \(0 < x^{*} < y\).
\end{exercise}

\begin{exercise}
Dada uma loteria \(L = (p_{1},\dots,p_{n})\) de possíveis resultados
\(x=(x_{1},\dots,x_{n})\), o equivalente certo \(C(L, u)\) é a
quantidade monetária que garante a utilidade média (i.e.~utilidade de
von Neumann-Morgenstern). Como a utilidade \(u(x) = - e^{-ax}\) é
definida pelo coeficiente \(a\), podemos escrever \(C(L, u) = C(L, a)\).

Portanto,

\begin{align*}
\cancel{ - }e^{-aC} &= \sum \limits_{i=1}^{n} p_{i} (\cancel{ - }e^{-ax_{i}}) \\
&= \sum \limits_{i=1}^{n} p_{i} e^{-ax_{i}}
\end{align*} Isolando \(C\), temos

\begin{align*}
C(L, a) &= -\frac{1}{a} \ln\left( \sum \limits_{i=1}^{n}p_{i} e^{-ax_{i}} \right) \\
&= \frac{1}{a}\ln \left( \frac{1}{\left( \sum \limits_{i=1}^{n}p_{i} e^{-ax_{i}} \right)} \right)
\end{align*} Defina-se, a fins de simplificação, variáveis auxiliares
\(\tilde{x}_{i}\) tais que
\[
x_{i} = \frac{1}{a} \ln \tilde{x}_{i} \iff \tilde{x}_{i}^{-1} = e^{-ax_{i}}
\]
Dessa forma, temos que \begin{align*}
C(L, a) &= \frac{1}{a}  \ln \left( \frac{1}{\left( \sum \limits_{i=1}^{n}p_{i} \frac{1}{\tilde{x}_{i}} \right)} \right) 
\end{align*}

Fazendo o mesmo para \(C\), i.e.~\(C = \frac{1}{a}\ln \tilde{C}\), temos
\begin{align*}
\tilde{C} &= \frac{1}{\left( \sum \limits_{i=1}^{n}p_{i} \frac{1}{\tilde{x}_{i}} \right)}  \\
&= \frac{\sum \limits_{i=1}^{n} p_{i}}{\left( \sum \limits_{i=1}^{n}p_{i} \frac{1}{\tilde{x}_{i}} \right)} 
\end{align*}

Ou seja, é possível escrever o equivalente certo de uma utilidade
exponencial como uma média \textit{harmônica ponderada}.

A fins de \textit{sanity check}: suponhamos que algum dos \(x_{i}\)'s
tende a \(0\)\footnote{Em verdade: consideremos \(x_{i}\) se aproximando
  de \(0\). Não queremos tomar o limite formal
  \(\lim\limits_{x\to_{0}}\) neste caso.} --- \textit{mas não sua
probabilidade} \(p_i\) ---, ou seja, é um caso que tende à perda máxima
do indivíduo com uma probabilidade não-nula.

Teremos, portanto, que, neste caso extremo,
\[
C \to \frac{1}{a}  \ln \left( \frac{x_{i}}{p_{i}}\right) 
\]

Ou seja, o equivalente certo torna-se (muito mais) dependente deste caso
extremo: o quanto se desejará pagar para evitar a aleatoriedade desta
loteria será contingente
\end{exercise}

\begin{exercise}
Com a utilidade \(u(x) = \frac{x^{1-\rho}-1}{1-\rho}\), temos que
\(C(L, u) = C(L, \rho)\), obtido da seguinte forma:

\begin{align*}
\frac{C^{1-\rho}-\cancel{ 1 }}{\cancel{ 1-\rho }} &= \sum \limits_{i=1}^{n} p_{i} \frac{x_{i}^{1-\rho}-\cancel{ 1 }}{\cancel{ 1-\rho }} \\
\end{align*}

Isolando \(C\), temos
\[
C(L, \rho) = \left( \sum \limits_{i=1}^{n} p_{i} x_{i}^{1-\rho} \right)^{1/(1-\rho)}
\]

Ou seja, o equivalente certo de uma utilidade isoelástica é uma
\href{https://en.wikipedia.org/wiki/Norm_(mathematics)\#p-norm}{p-norma}\footnote{É
  uma \(p\)-norma \textit{de facto} assumindo-se que \(x_{i}>0\).} ---
no caso, uma \((1-\rho)\)-norma --- \textbf{ponderada} (!!!).
\end{exercise}

\begin{exercise}
\begin{align*}
\ln C &= \sum \limits_{i=1}^{n} p_{i} \ln x_{i} \\
&= \sum \limits_{i=1}^{n} \ln x_{i}^{p_{i}} \\
&= \ln \left( \prod \limits_{i=1}^{n} x_{i}^{p_{i}} \right)
\end{align*}

Portanto,

\begin{align*}
C &= \prod \limits_{i=1}^{n} x_{i}^{p_{i}} \\
&= \left( \prod \limits_{i=1}^{n} x_{i}^{p_{i}} \right)^{1/\sum \limits_{i}p_{i}}
\end{align*}

Ou seja, o equivalente certo de uma utilidade logarítmica é uma média
geométrica \textit{ponderada}. Note-se que, aqui, não é uma média
\textit{aritmética}, ou seja, estamos falando de variações
\textit{multiplicativas} de riqueza, muito propício p.~ex. em contextos
financeiros.
\end{exercise}

\begin{exercise}
A probabilidade de que qualquer indivíduo \(i\) sofra uma perda
\(L_{i}\) é
\[
P(L_{i}) = p
\]

Supondo que haja \textit{um} indivíduo que sofra essa perda, a
quantidade que todos --- incluindo o próprio indivíduo --- pagariam para
ressarci-lo seria \(\frac{L}{N}\).

Como os riscos são independentes, a quantidade média de indivíduos que
será afetada é \(p\cdot N\). Multiplicando ambos, teremos a quantidade
que cada indivíduo deve pagar para ressarcir esta quantidade média de
indivíduos:
\[
q=\frac{L}{\cancel{ N }} \cancel{ N } p = pL
\]

Note-se que, embora a quantidade para ressarcir \textit{um único
indivíduo} seja inversamente proporcional a \(N\), a quantidade média de
\textit{indivíduos a se ressarcir} é diretamente proporcional a \(N\), o
que faz com que \(q\) não dependa de \(N\), e sim da probabilidade \(p\)
em si. Lembre-se que estamos considerando aqui \textbf{riscos
independentes}!!

Tomando um indivíduo indiferente ao risco com \(u(x)=ax\), pela
definição de equivalente certo, teremos \begin{align*}
u(w-q) = \cancel{ a }(w-q) &= (1-p)u(w) + p u(w-L) \\
&= (1-p) \cancel{ a }w + p \cancel{ a }w - p \cancel{ a }L \\
\implies q &= pL
\end{align*}
\end{exercise}

\begin{exercise}
Tomando as utilidades dos exercícios anteriores, teremos: 1- Exponencial

\begin{align*}
-e^{-a(w-q)} &= -(1-p)e^{-aw} - pe^{-a(w-L)} \\
\implies e^{-a(\cancel{ w }-q)} &= \cancel{ e^{-aw} }(1-p + p e^{aL}) \\
\implies q &= \frac{1}{a}\ln(1-p+pe^{aL})
\end{align*} 2 - Isoelástica (já cancelando os termos \(1-\rho\) do
denominador) \begin{align*}
(w-q)^{1-\rho}-\cancel{ 1 } &= (1-p)(w^{1-\rho}-\cancel{ 1 }) + p ((w-L)^{1-\rho}-\cancel{ 1 }) \\
\implies (w-q)^{1-\rho} &= w^{1-\rho}\left( 1-p + p \left( 1-\frac{L}{w} \right)^{1-\rho} \right) \\
\implies q &= w \left(1- \left( 1-p + p \left( 1-\frac{L}{w} \right)^{1-\rho} \right)\right)^{1/1-\rho}
\end{align*}

3 - Logarítmica \begin{align*}
\ln(w-q) &= (1-p)\ln w + p\ln(w-L) \\
&=\ln w^{1-p}+ \ln(w-L)^{p} \\
&= \ln(w^{1-p}(w-L)^{p})\\
\implies w-q &= w ^{1-p} (w-L)^{p} \\
&= w \left( 1-\frac{L}{w} \right)^{p}\\
\implies q = w\left( 1-\left( 1-\frac{L}{w} \right)^{p} \right)
\end{align*}

De novo: fique à vontade para testar a simulação
\href{https://www.desmos.com/calculator/quce5bk5v5}{Equivalentes certos
(Simulação Desmos)}!
\end{exercise}

\printbibliography

\end{document}
